% GENERATED FILE. Edit canonical lessons, then run npm run generate:books.
% canonical-source-hash: fnv1a64:0a5932a34d431724
% canonical-lessons: BN-C478-jotno, BN-C478-seba, BN-C478-tharmomitar, BN-C478-jolpotti, BN-C478-shushrusha-kora

\chapter{At the Sickbed}
\label{ch:bn-at-the-sickbed}
\hlchaptermodality{478}

\begin{chapteropening}
\textbf{By the end of this chapter:} \emph{I can talk about caring for someone who is ill.}

Complete the last lesson of chapter 478: I can talk about caring for someone who is ill.
\end{chapteropening}

\section[\texorpdfstring{\bn{যত্ন} (jôtnô)}{যত্ন (jôtnô)}]{\bn{যত্ন} (jôtnô) — care; looking after}

\label{lesson:BN-C478-jotno}



\begin{quote}
Before the new one: say the Bengali for to limp, then the Bengali for to moan.
\end{quote}

\subsection*{You'll want to know: \bn{যত্ন}}
\textbf{\bn{যত্ন}} — \emph{jôtnô} — \textquotedblleft{}care; looking after\textquotedblright{}.

\begin{quote}
\bn{নার্সরা} \bn{রোগীদের} \bn{যত্ন} \bn{নেন।}
\end{quote}

\begin{quote}
\emph{nārsrā rogīder jôtnô nen.}
\end{quote}

\begin{quote}
The nurses take care of the patients.
\end{quote}

\bn{যত্ন} \bn{নেওয়া}, to take care of; \bn{যত্ন} \bn{করে}, carefully.

\begin{cousinweb}
\textbf{\bn{যত্ন}} (\emph{jôtnô}) is a tatsama, Sanskrit \emph{yatna}, \textquotedblleft{}effort\textquotedblright{} (CDIAL 10404a); Bengali uses it for the effort of care.
\end{cousinweb}

\begin{grammarlens}[title={What you've built}]
One more word for describing what you lived and felt.
\end{grammarlens}

\subsection*{Your turn}
\begin{itemize}
\raggedright
  \item \emph{Say it:} \emph{jôtnô}
  \item \emph{Say it:} \emph{jôtnô}, once more
  \item \emph{Say it:} \emph{nārsrā rogīder jôtnô nen.} — The nurses take care of the patients.
  \item \emph{Recall:} say the Bengali for to limp, then the Bengali for to moan, then say \emph{jôtnô} again
\end{itemize}

\begin{tcolorbox}[breakable,colback=teal!4,colframe=teal!35!black,title={Before you move on}]
What does \bn{যত্ন} mean? (\textquotedblleft{}Care; looking after\textquotedblright{}.) Say it once more.
\end{tcolorbox}
\section[\texorpdfstring{\bn{সেবা} (sebā)}{সেবা (sebā)}]{\bn{সেবা} (sebā) — nursing; service}

\label{lesson:BN-C478-seba}



\begin{quote}
Before the new one: say the Bengali for to moan, then the Bengali for care.
\end{quote}

\subsection*{You'll want to know: \bn{সেবা}}
\textbf{\bn{সেবা}} — \emph{sebā} — \textquotedblleft{}nursing; service\textquotedblright{}.

\begin{quote}
\bn{মেয়ের} \bn{সেবায়} \bn{বাবা} \bn{তাড়াতাড়ি} \bn{সুস্থ} \bn{হলেন।}
\end{quote}

\begin{quote}
\emph{meyer sebāy bābā tāṛātāṛi susthô hôlen.}
\end{quote}

\begin{quote}
With his daughter's nursing, the father soon got well.
\end{quote}

\begin{cousinweb}
\textbf{\bn{সেবা}} (\emph{sebā}) is a tatsama, Sanskrit \emph{sevā}, \textquotedblleft{}attendance on, homage\textquotedblright{} (CDIAL 13595).
\end{cousinweb}

\begin{grammarlens}[title={What you've built}]
One more word for describing what you lived and felt.
\end{grammarlens}

\subsection*{Your turn}
\begin{itemize}
\raggedright
  \item \emph{Say it:} \emph{sebā}
  \item \emph{Say it:} \emph{sebā}, once more
  \item \emph{Say it:} \emph{meyer sebāy bābā tāṛātāṛi susthô hôlen.} — With his daughter's nursing, the father soon got well.
  \item \emph{Recall:} say the Bengali for to moan, then the Bengali for care, then say \emph{sebā} again
\end{itemize}

\begin{tcolorbox}[breakable,colback=teal!4,colframe=teal!35!black,title={Before you move on}]
What does \bn{সেবা} mean? (\textquotedblleft{}Nursing; service\textquotedblright{}.) Say it once more.
\end{tcolorbox}
\section[\texorpdfstring{\bn{থার্মোমিটার} (thārmomiṭār)}{থার্মোমিটার (thārmomiṭār)}]{\bn{থার্মোমিটার} (thārmomiṭār) — a thermometer}

\label{lesson:BN-C478-tharmomitar}



\begin{quote}
Before the new one: say the Bengali for care, then the Bengali for nursing.
\end{quote}

\subsection*{You'll want to know: \bn{থার্মোমিটার}}
\textbf{\bn{থার্মোমিটার}} — \emph{thārmomiṭār} — \textquotedblleft{}a thermometer\textquotedblright{}.

\begin{quote}
\bn{নার্সটি} \bn{থার্মোমিটার} \bn{দিয়ে} \bn{অরুণের} \bn{জ্বর} \bn{মাপলেন।}
\end{quote}

\begin{quote}
\emph{nārsṭi thārmomiṭār diye ôruner jbôr māplen.}
\end{quote}

\begin{quote}
The nurse took Arun's temperature with a thermometer.
\end{quote}

\begin{cousinweb}
\textbf{\bn{থার্মোমিটার}} (\emph{thārmomiṭār}) is borrowed from English \emph{thermometer}.
\end{cousinweb}

\begin{grammarlens}[title={What you've built}]
One more word for describing what you lived and felt.
\end{grammarlens}

\subsection*{Your turn}
\begin{itemize}
\raggedright
  \item \emph{Say it:} \emph{thārmomiṭār}
  \item \emph{Say it:} \emph{thārmomiṭār}, once more
  \item \emph{Say it:} \emph{nārsṭi thārmomiṭār diye ôruner jbôr māplen.} — The nurse took Arun's temperature with a thermometer.
  \item \emph{Recall:} say the Bengali for care, then the Bengali for nursing, then say \emph{thārmomiṭār} again
\end{itemize}

\begin{tcolorbox}[breakable,colback=teal!4,colframe=teal!35!black,title={Before you move on}]
What does \bn{থার্মোমিটার} mean? (\textquotedblleft{}A thermometer\textquotedblright{}.) Say it once more.
\end{tcolorbox}
\section[\texorpdfstring{\bn{জলপট্টি} (jôlpôṭṭi)}{জলপট্টি (jôlpôṭṭi)}]{\bn{জলপট্টি} (jôlpôṭṭi) — a cold compress (a wet cloth on the forehead)}

\label{lesson:BN-C478-jolpotti}



\begin{quote}
Before the new one: say the Bengali for nursing, then the Bengali for a thermometer.
\end{quote}

\subsection*{You'll want to know: \bn{জলপট্টি}}
\textbf{\bn{জলপট্টি}} — \emph{jôlpôṭṭi} — \textquotedblleft{}a cold compress (a wet cloth on the forehead)\textquotedblright{}.

\begin{quote}
\bn{জ্বর} \bn{কমাতে} \bn{দিদিমা} \bn{নাতির} \bn{কপালে} \bn{জলপট্টি} \bn{দিলেন।}
\end{quote}

\begin{quote}
\emph{jbôr kômāte didimā nātir kôpāle jôlpôṭṭi dilen.}
\end{quote}

\begin{quote}
To bring the fever down, the grandmother laid a cold compress on her grandson's forehead.
\end{quote}

\begin{cousinweb}
\textbf{\bn{জলপট্টি}} (\emph{jôlpôṭṭi}) puts \bn{জল} (taught), water, with \emph{pôṭṭi}, \textquotedblleft{}a strip of cloth\textquotedblright{}, probably from Sanskrit \emph{paṭṭa}, \textquotedblleft{}a slab\textquotedblright{} (CDIAL 7699), which Monier-Williams also gives for a bandage or strip of cloth.
\end{cousinweb}

\begin{grammarlens}[title={What you've built}]
One more word for describing what you lived and felt.
\end{grammarlens}

\subsection*{Your turn}
\begin{itemize}
\raggedright
  \item \emph{Say it:} \emph{jôlpôṭṭi}
  \item \emph{Say it:} \emph{jôlpôṭṭi}, once more
  \item \emph{Say it:} \emph{jbôr kômāte didimā nātir kôpāle jôlpôṭṭi dilen.} — To bring the fever down, the grandmother laid a cold compress on her grandson's forehead.
  \item \emph{Recall:} say the Bengali for nursing, then the Bengali for a thermometer, then say \emph{jôlpôṭṭi} again
\end{itemize}

\begin{tcolorbox}[breakable,colback=teal!4,colframe=teal!35!black,title={Before you move on}]
What does \bn{জলপট্টি} mean? (\textquotedblleft{}A cold compress (a wet cloth on the forehead)\textquotedblright{}.) Say it once more.
\end{tcolorbox}
\section[\texorpdfstring{\bn{শুশ্রূষা} \bn{করা} (shushrūshā kôrā)}{শুশ্রূষা করা (shushrūshā kôrā)}]{\bn{শুশ্রূষা} \bn{করা} (shushrūshā kôrā) — to nurse (a sick person)}

\label{lesson:BN-C478-shushrusha-kora}



\begin{quote}
Before the new one: say the Bengali for a thermometer, then the Bengali for a cold compress.
\end{quote}

\subsection*{You'll want to know: \bn{শুশ্রূষা} \bn{করা}}
\textbf{\bn{শুশ্রূষা} \bn{করা}} — \emph{shushrūshā kôrā} — \textquotedblleft{}to nurse (a sick person)\textquotedblright{}.

\begin{quote}
\bn{মীরা} \bn{রাত} \bn{জেগে} \bn{মামার} \bn{শুশ্রূষা} \bn{করল।}
\end{quote}

\begin{quote}
\emph{mīrā rāt jege māmār shushrūshā kôrlô.}
\end{quote}

\begin{quote}
Mira stayed up through the night nursing her uncle.
\end{quote}

Past (he/she …): \bn{তিনি} \bn{রোগীর} \bn{শুশ্রূষা} \bn{করলেন}, \emph{tini rogīr shushrūshā kôrlen}.

\begin{cousinweb}
\textbf{\bn{শুশ্রূষা} \bn{করা}} (\emph{shushrūshā kôrā}) is Sanskrit \emph{śuśrūṣā}, first \textquotedblleft{}the wish to hear\textquotedblright{} and then \textquotedblleft{}obedient service\textquotedblright{} (Monier-Williams), which Bengali uses for nursing the sick, with \emph{kôrā}, \textquotedblleft{}to do\textquotedblright{} (CDIAL 2814).
\end{cousinweb}

\begin{grammarlens}[title={What you've built}]
One more word for describing what you lived and felt.
\end{grammarlens}

\subsection*{Your turn}
\begin{itemize}
\raggedright
  \item \emph{Say it:} \emph{shushrūshā kôrā}
  \item \emph{Say it:} \emph{shushrūshā kôrā}, once more
  \item \emph{Say it:} \emph{mīrā rāt jege māmār shushrūshā kôrlô.} — Mira stayed up through the night nursing her uncle.
  \item \emph{Recall:} say the Bengali for a thermometer, then the Bengali for a cold compress, then say \emph{shushrūshā kôrā} again
\end{itemize}

\begin{tcolorbox}[breakable,colback=teal!4,colframe=teal!35!black,title={Before you move on}]
What does \bn{শুশ্রূষা} \bn{করা} mean? (\textquotedblleft{}To nurse (a sick person)\textquotedblright{}.) Say it once more.
\end{tcolorbox}
